\section{A rational experiment with explicit spectral constants}\label{sec:effective61}
The original rational pair in Theorem~\ref{thm:rationalseparation58}
now has a sharp linear fixed-error law, but its spectral constant is
still supplied existentially by Bourgain--Gamburd. We give a second
finite input for which every constant in a useful linear lower bound
is numerical. The deep input here is the classical
Lubotzky--Phillips--Sarnak theorem, not a new spectral-gap calculation.

Let $\sigma_1,\sigma_2,\sigma_3$ be the Pauli matrices in their usual
order, and put $D(x)=(I+\sum_jx_j\sigma_j)/2$ for $x\in\mathbb B_3$.
Thus $D(\mathbb B_3)=\mathcal D_2$ and
$\norm{D(x)-D(y)}_{\rm F}=\norm{x-y}/\sqrt2$.
Consider the following \emph{rational Bloch rotations}:
\begin{equation}\label{eq:lpsinput61}
 A_1=\frac15\begin{pmatrix}5&0&0\\0&-3&-4\\0&4&-3\end{pmatrix},\quad
 A_2=\frac15\begin{pmatrix}-3&0&4\\0&5&0\\-4&0&-3\end{pmatrix},\quad
 A_3=\frac15\begin{pmatrix}-3&-4&0\\4&-3&0\\0&0&5\end{pmatrix}.
\end{equation}
Take the alphabet
\[
 \mathcal A_{\rm LPS}=
       \{I,A_1,A_1^{\mathsf T},A_2,A_2^{\mathsf T},A_3,A_3^{\mathsf T}\}
\]
and the rational pure seed $e_1$, so $D(e_1)=P_0$ as in
\eqref{eq:rationaldata58}. The induced unitary lifts can be taken as
$(I-2i\sigma_j)/\sqrt5$. They are algebraic, not rational matrices;
the rational input in \eqref{eq:lpsinput61} describes the physical
Bloch action directly. This is a different alphabet from the rational
unitary pair of Theorem~\ref{thm:rationalseparation58}.

\begin{lemma}[The classical six-rotation spectral bound]\label{lem:lpsgap61}
Let $P_{\rm LPS}$ average uniformly over the six nonidentity letters
in $\mathcal A_{\rm LPS}$. On normalized $S^2$,
\begin{equation}\label{eq:lpsgap61}
 \norm{P_{\rm LPS}}_{L^2_0(S^2)\to L^2_0(S^2)}=\frac{\sqrt5}{3}.
\end{equation}
Their generated subgroup is dense in $SO(3)$.
\end{lemma}
\begin{proof}
The matrices in \eqref{eq:lpsinput61} are the adjoint actions of the
unit Hamilton quaternions
$(1+2\mathbf i)/\sqrt5,(1+2\mathbf j)/\sqrt5,(1+2\mathbf k)/\sqrt5$.
Together with their inverses they are exactly the norm-five set in the
Lubotzky--Phillips--Sarnak theorem \cite[Theorems~1.3 and~1.5]{LPS61};
see the precise restatement \cite[Theorem~1.1]{PLP61pre}, which gives
$2\sqrt p/(p+1)$ for this set at $p=5$. This proves
\eqref{eq:lpsgap61}. The cited theorem uses the theory of automorphic
forms and Deligne's theorem; no finite matrix enumeration replaces it.

Density can also be checked directly. Each rotation has cosine
$-3/5$. If its angle were a rational multiple of $2\pi$, twice this
cosine would be a rational algebraic integer, contrary to $-6/5\notin\Z$.
Its powers are dense in its circle subgroup. The three coordinate-axis
circle subgroups generate $SO(3)$, as follows either from their
infinitesimal generators and brackets or from Euler rotations. Hence
the closed generated subgroup is $SO(3)$.
\end{proof}

\begin{theorem}[An effective linear noisy law and an exponential endpoint]
\label{thm:effectivelps61}
For the fixed input \eqref{eq:lpsinput61}, the terminal target is
$D(A_w e_1)$, every decoder belongs to $\mathcal D_2$, and error is
Frobenius error. For $0<\varepsilon<1/\sqrt2$ put
$\zeta=1-\sqrt2\varepsilon$. Every admissible realization satisfies
\begin{equation}\label{eq:lpsoccupation61}
 \#\{1\le t\le N:|S_t|\le k\}
 \le 1+6\log(12k)+648\pi^2\frac{1-\zeta}{\zeta}k
 \le\frac{6500}{\zeta}k.
\end{equation}
In particular,
\begin{equation}\label{eq:lpslinear61}
 \frac{\zeta N}{6500}\le W_{N,\varepsilon}
                       \le52+\frac{104N}{\varepsilon}.
\end{equation}
For rational $\varepsilon$ in this interval the upper machine has
uniformly constructed rational legal decoders and dyadic stochastic
rows, each supported on at most four labels. It uses at most $Nb$
fair bits for hidden transitions, where
$b=\lceil\log_2(24(N+1)/\varepsilon)\rceil$.

At every horizon $N\ge0$, the simultaneous exact minimum profile is
\begin{equation}\label{eq:lpsexact61}
                  |S_t|_{\min}=5^t\quad(0\le t\le N).
\end{equation}
The same profile remains necessary and is attained throughout
\begin{equation}\label{eq:lpsrobust61}
                         0\le\varepsilon\le25^{-N}/16.
\end{equation}
For $\varepsilon\ge1/\sqrt2$, one label suffices.
\end{theorem}
\begin{proof}
For a spherical chordal cap of radius $h\le1$, normalized area is
exactly $h^2/4$. Consequently the support of a mixture of $k$ kernels
$K_h$ occupies at most $kh^2/4$ of each group orbit. Take $h=k^{-1/2}$.
In Lemma~\ref{lem:entropydefect61} we may take
\[
 \theta=1/4,\qquad \kappa^2=5/9,\qquad
 \gamma=(1-\kappa^2)(1-\theta)=1/3.
\]
Lemma~\ref{lem:smoothentropy61} gives $L_h=\log(12k)$ and $B_3=24$.
The proof of Theorem~\ref{thm:sharpoccupation61}, with these constants
and the normalization $\zeta=1-\sqrt2\varepsilon$, gives the first
inequality in \eqref{eq:lpsoccupation61} as a real bound before taking
an integer ceiling. For $k\ge1$, $\log(12k)\le3k$, and $\pi^2<10$.
Thus its right side is at most
$19k+6480(1-\zeta)k/\zeta\le6500k/\zeta$.
Counting all positive cuts proves the lower bound in
\eqref{eq:lpslinear61}.

For the upper bound apply the rational compiler of
Theorem~\ref{thm:rationalcompiler58} to the displayed seven-letter list
and the first-axis seed. At rational tolerance take
$\beta=1-\varepsilon/2$ and
$m=\max\{2,\lceil\sqrt{2N/\varepsilon}\rceil\}$.
It gives an exact rational realization of $D(\beta A_we_1)$ on at most
$13m^2\le52+52N/\varepsilon$ labels, with sparse common rows.
The target change has Frobenius norm $\varepsilon/(2\sqrt2)$.
The stated dyadic precision adds at most
$3\sqrt2N2^{-b}\le\sqrt2\varepsilon/8$, so the total is less than
$\varepsilon$. This proves the looser advertised numerical upper bound
and the fair-bit statement. For a real tolerance choose a rational
$\widetilde\varepsilon\in(\varepsilon/2,\varepsilon)$ and use the
same construction at $\widetilde\varepsilon$; the bound
$52+104N/\varepsilon$ follows. The algorithmic assertion uses rational
input, not an oracle for an unencoded real number.

We next prove the exact profile independently of the spectral bound.
Lemma~\ref{lem:freequaternion58} shows that the three norm-five
quaternion generators freely generate modulo scalars, hence their
adjoint rotations form a free group on $a,b,c$. The stabilizer of
$e_1$ in this free group is exactly $\langle a\rangle$. Indeed, a
rotation fixing $e_1$ commutes with the rotation $a$ about that axis.
In a free group the centralizer of this generator is its cyclic
subgroup: write a reduced word outside $\langle a\rangle$ as
$a^rva^s$, where $v$ is nonempty and begins and ends with neither
$a$ nor $a^{-1}$. Its conjugate of $a$ reduces to
$a^r v a v^{-1}a^{-r}$, which is not $a$. The converse, that every
power of $a$ fixes $e_1$, is immediate.

The orbit is therefore indexed by the right cosets of $\langle a\rangle$.
Each coset has a unique reduced representative with no terminal
$a^{\pm1}$, obtained by removing the maximal terminal power of $a$.
It has least word length in that coset. At length $\ell\ge1$ there
are $4\cdot5^{\ell-1}$ such representatives. Identity padding gives
exactly the ball of these cosets at command length $t$, whose cardinality
is $1+4\sum_{\ell=1}^t5^{\ell-1}=5^t$.
Theorem~\ref{thm:extreme57} applies because pure states are extreme in
$\mathcal D_2$, proving \eqref{eq:lpsexact61} at every cut; deterministic
orbit storage attains the entire profile at once.

For robustness, distinct cut-$t$ Bloch targets have coordinates in
$5^{-t}\Z^3$ and hence distance at least $5^{-t}$. One fixed suffix
of length $N-t$ preserves these distances. If $z_i\in\mathbb B_3$
is the conditional terminal decoder mean from cut label $i$,
wordwise error at a pure target $n$ implies
\[
 \sum_i p_i\norm{z_i-n}^2\le2\sqrt2\varepsilon.
\]
Some positive-mass label is within $\sqrt{2\sqrt2\varepsilon}$ of $n$.
Under \eqref{eq:lpsrobust61}, twice this radius is strictly less than
$5^{-N}\le5^{-t}$. The corresponding closed balls for distinct
prefix orbit points are disjoint, and so require distinct labels.
The exact construction attains this bound. No fixed positive tolerance
is asserted to lie in \eqref{eq:lpsrobust61} for all horizons.
Finally $D(0)=I/2$ is at Frobenius distance $1/\sqrt2$ from every pure
state, proving the one-label assertion.
\end{proof}

\begin{remark}[What is effective]\label{rem:effective61}
The input, orbit counts, robust tolerance, spectral norm, constants in
\eqref{eq:lpsoccupation61}--\eqref{eq:lpslinear61}, and rational compiler
are explicit. The constant $6500$ is conservative; it is not an optimal
leading constant. The spectral norm is imported from
\cite{LPS61,PLP61pre}, not inferred from finite harmonic tests.
For the earlier five-letter rational-unitary example,
Theorem~\ref{thm:sharpoccupation61} removes the logarithm but does not
supply a numerical Bourgain--Gamburd constant. Neither statement claims
a complete two-parameter phase diagram as the error tends to zero.
The finite-table and fair-bit resources refer to a classical hidden
label and a numerical matrix output, not quantum-state preparation.
\end{remark}
